Primary screening reduces coordinate work from $O(n_cd)$ to $O(cd)$,
but each survivor still requires a full residual-distance calculation.
Sharing one candidate among several lanes shortens each lane's loop.
The budget stays fixed while this section changes how that work is
distributed.

\subsection{Reconstructed scoring and the query-only constant}

The stored correction $\gamma=2\langle\widehat p,\widehat r\rangle$
combines primary and residual contributions into the refined score.
Expanding in real arithmetic explains what this score ranks:
\begin{align}
 \refineddist(q,x)
 &=\primarydist(q,x)+\|q'-\widehat r(x)\|_2^2+\gamma(x)
 \label{eq:5-refined-score}\\
 &=2\|q'\|_2^2+\|\widehat p+\widehat r\|_2^2-2\langle q',\widehat p+\widehat r\rangle\nonumber\\
 &=\|q'-\widehat p(x)-\widehat r(x)\|_2^2+\|q'\|_2^2.
 \label{eq:5-refined-constant}
\end{align}
The added query norm is the same for every candidate, so the ranking is
that of reconstructed squared distance \cite{jhq}. Scalar execution
sums residual squared differences before adding the primary score and
correction. The query and any codebooks that fit are staged in shared
memory.

\subsection{Mapping lanes to words and reducing partial sums}

Let $\lrefine=L\in\{1,8,16,32\}$ be the lanes assigned to one candidate.
Each lane visits a strided sequence of residual-code words:
\[
 d_{\rm word}=32/B_R,\qquad
 \text{words of lane }a:\ a,a+L,a+2L,\ldots.
\]
A word holds four eight-bit codes or eight four-bit codes. Lanes sum
their own decoded contributions, then exchange partial sums using
butterfly shuffles, a sequence of pairwise lane exchanges. Subgroup lane
zero adds the primary score and correction and writes the result.

\par\medskip\noindent\begin{minipage}{\linewidth}
\captionsetup{hypcap=false}
\centering
\begin{tikzpicture}[font=\small\linespread{1}\selectfont,
 box/.style={draw=ThesisNavy!60,rounded corners=2pt,align=center,text width=6.0cm,minimum height=1.25cm,inner sep=4pt},
 arr/.style={-{Latex[length=2mm]},draw=ThesisNavy!80,thick}]
\node[box,fill=ThesisNavy!8] (a) at (0,0) {Warp lanes 0--15: candidate $x_0$\\Local lane $a$: words $a,a+16,\ldots$};
\node[box,fill=ThesisNavy!8] (b) at (7.3,0) {Warp lanes 16--31: candidate $x_1$\\Local lane $a$: words $a,a+16,\ldots$};
\node[box] (ar) at (0,-1.9) {Decode and accumulate\\Sum within lanes 0--15};
\node[box] (br) at (7.3,-1.9) {Decode and accumulate\\Sum within lanes 16--31};
\node[box,fill=ThesisNavy!3] (outa) at (0,-3.8) {Lane 0 writes $D_H(q,x_0)$\\Add own primary score and correction};
\node[box,fill=ThesisNavy!3] (outb) at (7.3,-3.8) {Lane 16 writes $D_H(q,x_1)$\\Add own primary score and correction};
\foreach \x/\y in {a/ar,b/br,ar/outa,br/outb}\draw[arr] (\x)--(\y);
\end{tikzpicture}

\captionof{figure}[Sixteen-lane residual refinement]{Sixteen-lane residual refinement.}
\label{fig:5-lane-refinement}
\end{minipage}\par\medskip
With sixteen lanes, a warp processes two candidates independently, as
shown in \cref{fig:5-lane-refinement}. Local lane $a$ visits words
$a,a+16,\ldots$; warp lanes 0 and 16 write the separate outputs.
Word-aligned rows and $d_{\rm word}\mid D_s$ keep each word within its
codebook. All lanes participate in shuffles; invalid slots contribute
zero and ultimately receive an infinite output score.

The experimental names describe this mapping: L0 reads bytes with one
lane, L1 reads words with one lane, and L8/L16/L32 use that many lanes
per candidate. More lanes shorten local loops but add shuffle work and
reduce the candidates processed concurrently. \Cref{sec:7-3} compares
these choices and their numerical effects.

\subsection{Final selection and the numerical boundary}

Final selection scans all retained slots in each output round. Assuming
at least $k$ finite valid survivors with unique IDs, it chooses
\[
 i^*=\arg\min_{i<c}(D_{H,i},\mathrm{ID}_i),\quad d^*=D_{H,i^*},\quad
 D_{H,i^*}\leftarrow+\infty
\]
and emits $(\mathrm{ID}_{i^*},d^*)$. Replacing that score by infinity
makes the next round select the next candidate, yielding the exact
retained-pool top-$k$ without per-thread quotas. With too few finite
entries, the kernel can repeat the smallest valid ID at infinity;
all-invalid pools return invalid IDs. The evaluated settings satisfy
the finite-entry premise.

Primary ties use physical position and final ties use database ID.
This final selector is exact for its input scores, while routing and
quantised reconstruction still determine raw-vector recall. Changing
the summation order can change FP32 rounding and reorder near ties.
